Durante el período comprendido entre 2015 y 2025, los principales agregados monetarios de Guatemala muestran una tendencia general de crecimiento sostenido. El Producto Interno Bruto (PIB) presenta la evolución más estable y moderada, reflejando un crecimiento económico constante a lo largo del período. Por su parte, los medios de pago (M2) y el crédito bancario al sector privado registran un crecimiento superior al del PIB, manteniendo igualmente una trayectoria estable y con baja volatilidad. Finalmente, el numerario en circulación y las remesas familiares son las variables que experimentan el mayor crecimiento relativo, aunque las remesas destacan por presentar fluctuaciones más pronunciadas durante algunos trimestres.
Las diferencias en el comportamiento de estas variables pueden asociarse a su naturaleza económica. Mientras que el PIB, el crédito bancario, los medios de pago y el numerario en circulación responden principalmente a las condiciones macroeconómicas internas y a la política monetaria implementada por el Banco de Guatemala, las remesas familiares dependen en mayor medida de factores externos, como la situación económica de los países de origen de los migrantes, las condiciones del mercado laboral internacional y los cambios en las políticas migratorias. Esta dependencia explica la mayor volatilidad observada en dicha variable.
Desde una perspectiva de política económica, los resultados evidencian que el seguimiento de las remesas familiares y del numerario en circulación resulta especialmente relevante debido a la magnitud de su crecimiento y a la influencia que ejercen sobre la liquidez de la economía. No obstante, el análisis de estos indicadores debe complementarse con variables como el crédito bancario, los medios de pago y el PIB, ya que en conjunto proporcionan una visión más completa de la evolución del sistema monetario y financiero del país y permiten fundamentar de mejor manera la toma de decisiones.